\documentclass[12pt]{article}

% Packages
\usepackage[margin=2.5cm]{geometry}
\usepackage{amsmath}
\usepackage{amssymb}
\usepackage{booktabs}
\usepackage{enumitem}
\usepackage[hidelinks]{hyperref}

\setlength{\parindent}{0pt}
\setlength{\parskip}{6pt}

\begin{document}

\title{Evaluation of ``A Universal Route to $Q_{\max}$ and $K_{\mathrm{aff}}$\\
for PFAS Adsorption in Solids''}
\author{Evaluation Report}
\date{\today}
\maketitle

% ======================================================================
\section{Executive summary}
% ======================================================================

The manuscript proposes a model-agnostic workflow for extracting the
saturation capacity $Q_{\max}$ and a capacity-weighted mean affinity
constant $K_{\mathrm{aff}}$ from adsorption data, using the heterogeneous
Langmuir integral as its theoretical backbone.  An optional partition
term $K_p C$ is included for PFAS systems.  A nonparametric distribution
recovery via regularized non-negative least squares (NNLS) is offered as
a secondary tool.

\textbf{Overall assessment:} The manuscript is a competent synthesis of
well-established theoretical results, repackaged as a practical workflow.
The core mathematical content---the heterogeneous Langmuir integral, the
moment relations for $Q_{\max}$ and $K_H$, and the NNLS-based
distribution recovery---are all published results dating back decades.
The contribution is therefore primarily methodological/pedagogical rather
than theoretically novel.  Several technical issues require correction.

% ======================================================================
\section{Detailed originality analysis}
% ======================================================================

\subsection{Element 1: The heterogeneous Langmuir integral}

\textbf{Manuscript claim:} Eq.~(1) represents adsorption as
$q(C)=\int_0^\infty \rho(K)\,KC/(1+KC)\,dK$.

\textbf{Prior art:} This is the standard integral equation of adsorption
on heterogeneous surfaces, established since at least the 1940s:
\begin{itemize}[nosep]
  \item Sips, R.\ (1948). ``On the structure of a catalyst surface,''
        \textit{J.\ Chem.\ Phys.}\ \textbf{16}, 490--495.  Derived the
        Sips (Langmuir--Freundlich) isotherm by solving this integral
        with a quasi-Gaussian (logistic) distribution.
  \item Sips, R.\ (1950). ``On the structure of a catalyst surface.\ II,''
        \textit{J.\ Chem.\ Phys.}\ \textbf{18}, 1024--1026.
  \item T\'oth, J.\ (1971). ``State equations of the solid-gas interface
        layers,'' \textit{Acta Chim.\ Acad.\ Sci.\ Hung.}\ \textbf{69},
        311--328.
  \item Jaroniec, M.\ and Madey, R.\ (1988). \textit{Physical Adsorption
        on Heterogeneous Solids}, Elsevier, Amsterdam.  Comprehensive
        treatment of the integral equation, energy distributions, and
        solution methods.
  \item Rudzinski, W.\ and Everett, D.H.\ (1992). \textit{Adsorption of
        Gases on Heterogeneous Surfaces}, Academic Press, London.
        Definitive monograph covering analytical and numerical inversion
        of the integral equation.
\end{itemize}

\textbf{Verdict:} \textbf{Not novel.}  Textbook material for at least
35~years.

\subsection{Element 2: $Q_{\max}$ as the zeroth moment of $\rho(K)$}

\textbf{Manuscript claim:} Eq.~(5): $Q_{\max}=\int_0^\infty\rho(K)\,dK$
and $\lim_{C\to\infty}q(C)=Q_{\max}$.

\textbf{Prior art:} This is a direct consequence of dominated convergence
applied to the Langmuir kernel $KC/(1+KC)\to 1$ as $C\to\infty$.  It
appears in every treatment of heterogeneous adsorption (Jaroniec \&
Madey 1988, Rudzinski \& Everett 1992, and many others).

\textbf{Verdict:} \textbf{Not novel.}

\subsection{Element 3: $K_H$ as the first moment and $K_{\mathrm{aff}}=K_H/Q_{\max}$}

\textbf{Manuscript claim:} Eqs.~(7)--(9): $K_H=Q_{\max}\int K f(K)\,dK$
and $K_{\mathrm{aff}}=\int K f(K)\,dK$ is the capacity-weighted mean
affinity.

\textbf{Prior art:}
\begin{itemize}[nosep]
  \item The Henry constant as the first moment of $\rho(K)$ is a standard
        result (Jaroniec \& Madey 1988, Chapter~3; Rudzinski \& Everett
        1992, Chapter~5).
  \item The ratio $K_H/Q_{\max}$ as a mean affinity is used implicitly
        in binding theory since at least Klotz, I.M.\ and Hunston, D.L.\
        (1971), ``Properties of graphical representations of multiple
        classes of binding sites,'' \textit{Biochemistry} \textbf{10},
        3065--3069.
  \item Nederlof, M.M., van Riemsdijk, W.H., and Koopal, L.K.\ (1990).
        ``Determination of adsorption affinity distributions,''
        \textit{J.\ Colloid Interface Sci.}\ \textbf{135}, 410--426.
        General framework for affinity distribution recovery.
  \item Rampey, A.M., Umpleby, R.J., Rushton, G.T., Iseman, J.C.,
        Shah, R.N., and Shimizu, K.D.\ (2004).  ``Characterization of
        the imprint effect and the influence of imprinting conditions on
        affinity, capacity, and heterogeneity in molecularly imprinted
        polymers using the Freundlich isotherm-affinity distribution
        analysis,'' \textit{Anal.\ Chem.}\ \textbf{76}, 1123--1133.
        \emph{Explicitly extract $Q_{\max}$-like parameters and affinity
        distributions from the heterogeneous Langmuir integral.}
\end{itemize}

\textbf{Verdict:} \textbf{Not novel.}  The notation $K_{\mathrm{aff}}$ may
be new but the concept is not.

\subsection{Element 4: Partition-plus-Langmuir dual-mode model}

\textbf{Manuscript claim:} Eq.~(4): $q(C)=K_p C+\int\rho(K)\,KC/(1+KC)\,dK$.

\textbf{Prior art:}
\begin{itemize}[nosep]
  \item Vieth, W.R., Howell, J.M., and Hsieh, J.H.\ (1976). ``Dual
        sorption theory,'' \textit{J.\ Membr.\ Sci.}\ \textbf{1}, 177--220.
        The dual-mode sorption model $q=k_D C + C'_H bC/(1+bC)$ has been
        standard in polymer membrane science for 50~years.
  \item Paul, D.R.\ and Koros, W.J.\ (1976). ``Effect of partially
        immobilizing sorption on permeability and the diffusion time lag,''
        \textit{J.\ Polym.\ Sci.}\ \textbf{14}, 675--685.
  \item For PFAS specifically, numerous recent papers apply dual-mode
        (partition + adsorption) models to soils and sediments.
\end{itemize}

\textbf{Verdict:} \textbf{Not novel.}  Direct application of the 1976
dual-mode sorption model.

\subsection{Element 5: $1/C$ extrapolation for $Q_{\max}$}

\textbf{Manuscript claim:} Eq.~(12): $q(C)\approx Q_{\max}-A_1/C$ for
large $C$, used to extrapolate $Q_{\max}$.

\textbf{Prior art:} This is the high-concentration asymptotic expansion
of the heterogeneous Langmuir integral.  The $1/C$ correction term
appears in:
\begin{itemize}[nosep]
  \item Jaroniec, M.\ (1983). ``Physical adsorption on heterogeneous
        solids,'' \textit{Adv.\ Colloid Interface Sci.}\ \textbf{18},
        149--225.
  \item The $q$ vs.\ $1/C$ linear extrapolation is conceptually related
        to the classical Langmuir linearization ($1/q$ vs.\ $1/C$), which
        dates to Langmuir (1918) and Lineweaver--Burk (1934) in enzyme
        kinetics.
\end{itemize}

\textbf{Verdict:} \textbf{Not novel}, though the specific use for
heterogeneous systems is a reasonable (if straightforward) extension.

\subsection{Element 6: Nonparametric distribution recovery via regularized NNLS}

\textbf{Manuscript claim:} Eqs.~(18)--(20): Discretize the integral into a matrix equation $q\approx A\pm{q}_m$ and solve with Tikhonov-regularized NNLS.

\textbf{Prior art:}
\begin{itemize}[nosep]
  \item Umpleby, R.J., Baxter, S.C., Chen, Y., Shah, R.N., and
        Shimizu, K.D.\ (2001). ``Characterization of molecularly imprinted
        polymers with the Langmuir--Freundlich isotherm,''
        \textit{Anal.\ Chem.}\ \textbf{73}, 4584--4591.
  \item Umpleby, R.J., Baxter, S.C., Rampey, A.M., Rushton, G.T.,
        Chen, Y., and Shimizu, K.D.\ (2004). ``Characterization of the
        heterogeneous binding site affinity distributions in molecularly
        imprinted polymers,'' \textit{J.\ Chromatogr.\ B} \textbf{804},
        141--149.  Used discrete Langmuir mixtures with numerical
        distribution recovery.
  \item Szabelski, P., Kierzek, A.M., et al.\ (2002).  ``Determination of
        single component isotherms and affinity energy distributions by
        chromatography,'' \textit{J.\ Chromatogr.\ A} \textbf{988}, 41--55.
  \item Cernik, M., Borkovec, M., and Westall, J.C.\ (1995).
        ``Regularized least-squares methods for the calculation of discrete
        and continuous affinity distributions for heterogeneous sorbents,''
        \textit{Environ.\ Sci.\ Technol.}\ \textbf{29}, 413--425.
        \emph{Directly used regularized NNLS for affinity distributions
        in environmental sorbents---essentially the same approach as the
        manuscript under review.}
  \item Carter, M.C.\ and Kilduff, J.E.\ (1995).
        ``Site energy distribution analysis of preloaded adsorbents,''
        \textit{Environ.\ Sci.\ Technol.}\ \textbf{29}, 1773--1780.
        Site energy distribution analysis for activated carbon using
        numerical inversion.
  \item In the broader inverse problems literature, regularized NNLS for
        Fredholm integral equations of the first kind is a standard
        technique (Tikhonov \& Arsenin, 1977; Hansen, 1998).
\end{itemize}

\textbf{Verdict:} \textbf{Not novel.}  Well-established technique in both
adsorption science and inverse problems.

\subsection{Element 7: Bootstrap uncertainty estimation}

\textbf{Manuscript claim:} Nonparametric bootstrap for confidence
intervals on $Q_{\max}$ and $K_{\mathrm{aff}}$.

\textbf{Prior art:} Bootstrap is a standard statistical method (Efron,
1979).  Its application to adsorption parameters is routine.

\textbf{Verdict:} \textbf{Not novel.}

% ======================================================================
\section{Summary of originality assessment}
% ======================================================================

\begin{table}[h!]
\centering
\caption{Originality of each element in the manuscript.}
\begin{tabular}{@{}lll@{}}
\toprule
\textbf{Element} & \textbf{Novel?} & \textbf{Key prior reference} \\
\midrule
Heterogeneous Langmuir integral & No & Sips (1948); Jaroniec \& Madey (1988) \\
$Q_{\max}$ as zeroth moment     & No & Jaroniec \& Madey (1988) \\
$K_H$ as first moment           & No & Jaroniec \& Madey (1988) \\
$K_{\mathrm{aff}}=K_H/Q_{\max}$ (mean affinity)  & No & Klotz \& Hunston (1971) \\
Dual-mode (partition + Langmuir) & No & Vieth et al.\ (1976) \\
$1/C$ extrapolation for $Q_{\max}$ & No & Jaroniec (1983); Langmuir (1918) \\
NNLS distribution recovery       & No & Umpleby et al.\ (2001, 2004) \\
Bootstrap uncertainty             & No & Efron (1979) \\
\midrule
\textit{Packaging as a workflow} & \textit{Partial} &
  \textit{Synthesis of known methods} \\
\textit{PFAS-specific application} & \textit{Partial} &
  \textit{Application context is timely} \\
\bottomrule
\end{tabular}
\end{table}

% ======================================================================
\section{Technical issues requiring correction}
% ======================================================================

\subsection{Issue 1: Convergence of $K_H$ is not guaranteed}

The manuscript states Eq.~(7): $K_H=\int_0^\infty K\rho(K)\,dK$ without
specifying when this integral converges.  For important distribution
classes, it does not:

\begin{itemize}[nosep]
  \item \textbf{Freundlich:} The site-energy distribution is a power law
        $f(E)\propto \exp(nE/RT)$, which in $K$-space gives
        $f(K)\propto K^{n-1}$.  The first moment
        $\int K f(K)\,dK \propto \int K^n\,dK$ diverges for $n<1$
        (the physically relevant regime).
  \item \textbf{Sips with $m<1$:} The logistic distribution similarly has
        a divergent first moment.
\end{itemize}

The manuscript should state explicitly that $K_{\mathrm{aff}}$ is
well-defined only when $\int K\rho(K)\,dK<\infty$, and discuss what
happens for power-law-tailed distributions (where Henry's law does not
hold in the strict mathematical sense, and $K_H$ must be defined
operationally over a finite concentration window).

\subsection{Issue 2: The $1/C$ extrapolation requires $A_1<\infty$}

Eq.~(12) assumes $A_1=\int_0^\infty \rho(K)/K\,dK<\infty$.  For
distributions with substantial weight near $K=0$ (e.g., some Toth
distributions with strong asymmetry), this integral may diverge.  The
manuscript should state this condition and provide a diagnostic.

\subsection{Issue 3: Iterative $K_p$ estimation lacks convergence analysis}

The manuscript suggests iterating Steps~2 and~3 ``once or twice to
self-consistency'' without proving convergence or providing a stopping
criterion.  A proper alternating estimation procedure should include:
\begin{itemize}[nosep]
  \item A convergence criterion (e.g., relative change in
        $\widehat{K}_p$ below a tolerance).
  \item Discussion of uniqueness: the decomposition
        $q=K_p C + q_{\mathrm{sat}}$ is not identifiable when the
        saturating component has not reached its plateau.
\end{itemize}

\subsection{Issue 4: Selection of $\mathcal{I}_0$ and $\mathcal{I}_\infty$ subsets}

The manuscript provides only vague guidance (``increase the number of
included points until\ldots'').  The corrected version should include:
\begin{itemize}[nosep]
  \item A quantitative criterion for Henry region selection (e.g., based
        on the curvature $q''(C)$ or an $F$-test comparing linear vs.\
        quadratic fits).
  \item A quantitative criterion for plateau region selection.
  \item Sensitivity analysis showing how $K_{\mathrm{aff}}$ varies with
        the choice of subsets.
\end{itemize}

\subsection{Issue 5: Regularization parameter selection}

The NNLS formulation (Eq.~20) includes $\lambda$ but provides no
guidance for choosing it.  Standard methods should be mentioned:
\begin{itemize}[nosep]
  \item L-curve method (Hansen, 1992).
  \item Generalized cross-validation (GCV).
  \item Discrepancy principle (Morozov, 1966).
\end{itemize}

\subsection{Issue 6: Missing literature citations}

The manuscript contains \emph{zero} references.  A publication-ready
version must cite the foundational works listed in Section~2.

\subsection{Issue 7: Assumption A1 is restrictive}

Assumption A1 (independent-site superposition, no lateral interactions or
cooperativity) excludes:
\begin{itemize}[nosep]
  \item Systems with cooperative binding (Hill isotherm, sigmoidal
        isotherms).
  \item Systems with adsorbate--adsorbate interactions on the surface
        (Fowler--Guggenheim, Temkin deviation from ideality).
  \item Multilayer adsorption (BET-type).
\end{itemize}
This should be discussed more explicitly.

% ======================================================================
\section{Recommendations}
% ======================================================================

\begin{enumerate}
  \item \textbf{Reframe the contribution:} The manuscript should be
        presented as a practical workflow and decision protocol, not as
        a theoretically novel result.  The title ``Universal Route'' is
        overstated; something like ``A Practical Model-Agnostic Protocol
        for Extracting $Q_{\max}$ and $K_{\mathrm{aff}}$'' would be more
        accurate.
  \item \textbf{Add comprehensive citations} to the foundational
        literature (Sips, Jaroniec, Rudzinski, Umpleby, Shimizu, Vieth,
        Nederlof, etc.).
  \item \textbf{State convergence conditions} for $K_H$ and $A_1$
        explicitly.
  \item \textbf{Formalize the iterative procedure} for $K_p$ with a
        convergence criterion.
  \item \textbf{Provide quantitative criteria} for selecting Henry and
        plateau regions.
  \item \textbf{Include regularization parameter selection} guidance for
        the NNLS step.
  \item \textbf{Add validation:} Test on synthetic data generated from
        known distributions (Sips, Toth, Double Langmuir) and on
        published PFAS datasets.
  \item \textbf{Add diagnostic plots:} residual analysis, sensitivity
        of $K_{\mathrm{aff}}$ to subset choices, comparison with
        parametric model fits.
\end{enumerate}

% ======================================================================
\section{Conclusion}
% ======================================================================

The manuscript assembles known theoretical results into a practical
workflow.  Its value lies in the synthesis and in targeting the PFAS
community, which may not be familiar with the heterogeneous adsorption
literature from physical chemistry and surface science.  However, the
manuscript currently (a)~does not cite any prior work, (b)~contains
technical gaps regarding convergence conditions and parameter selection,
and (c)~overstates the novelty of the approach.  A revised version
addressing these issues (provided as \texttt{Universal\_route\_v2.tex})
would be a useful methodological contribution.

\end{document}
